\section{Nonvanishing from signed phases}\label{sec:B}

We seek sections of a bounded positive multiple of a rational
polarization, even when its denominators are unbounded.  The hypothesis
will control the fractional orders of the polarization at primes whose
discrepancies tend to zero.  The sign of that control will determine
which corrections can be removed in the proof.

For a rational Weil divisor $H$ on a normal variety, set
\begin{equation}\label{eq:B-ramification-boundary}
 R(H)=\sum_P\left(1-\frac1{r_P}\right)P,
 \qquad
 r_P=\min\{r\in\Z_{>0}:r\operatorname{coeff}_P(H)\in\Z\}.
\end{equation}
Only finitely many terms are nonzero.  When $H$ is $\Q$-Cartier, its
order at a higher divisor means the order of its Cartier b-pullback,
rather than a rounded trace.  A \emph{signed phase} of this order is a lift of its image in
$\R/\Z$; whenever this lift tends to zero, it is eventually the unique
lift in $(-1/2,1/2)$.

We use the comparison results of Section~\ref{sec:A} and two boundedness
theorems.  Rationally connected projective varieties carrying effective
generalized $\delta$-lc Calabi--Yau data in bounded dimension are bounded
up to isomorphism in codimension one
\cite[Theorem~1.7]{BirkarCY}.  No bound on the Cartier index of the nef
b-divisor occurs in that theorem.  We also use effective birationality
for a nef and big integral Weil divisor $N$ on a projective
$\delta$-lc variety when $N-K$ is pseudo-effective
\cite[Theorem~1.1]{BirkarPolarised}.  Our applications of the latter
theorem take place on $\Q$-factorial models, so that $N$ is
$\Q$-Cartier, but need not be Cartier.

\begin{lemma}[Phase nonvanishing]\label{lem:B-phase-nonvanishing}
Let $(V_i,H_i)$ be a sequence in bounded dimension, where $V_i$ is
projective of Fano type and $H_i$ is a nef and big $\Q$-Cartier rational
Weil divisor.  Suppose that, for rational $t\in[0,L_i]$, where
$L_i\to\infty$, there are generalized data with discrepancies $a_i(t)$
such that
\begin{equation}\label{eq:B-family}
 B_i(t)\ge R(H_i),\qquad
 K_{V_i}+B_i(t)+\mathbf M_i(t)_{V_i}\sim_{\Q}tH_i,
 \qquad \mathbf M_i(t)\text{ is b-nef and $\Q$-Cartier}.
\end{equation}
Assume that $a_i(0)$ is klt and that $a_i(t)$ is nondecreasing in $t$
at every divisorial valuation.

Assume the following condition, also after restriction to any
subsequence.  For every choice of rational $t_i\to\infty$ in the
parameter range, one can pass to a further subsequence and choose a
fixed positive integer $q$ such that, for every sequence of primitive
divisorial valuations $T_i$ with $a_i(t_i,T_i)\to0$, the signed phases
$\theta_i(T_i)$ of $qH_i(T_i)$ satisfy
\begin{equation}\label{eq:B-phase-hypothesis}
 -o\bigl(a_i(t_i,T_i)\bigr)
 \le \theta_i(T_i)
 \le O\bigl(a_i(t_i,T_i)\bigr).
\end{equation}
Then there is an integer $M>0$ such that for every $i$ some
$m_i\in\{1,\ldots,M\}$ satisfies
$|\lfloor m_iH_i\rfloor|\ne\varnothing$.
Equivalently, after replacing $M$ by a common multiple, the same
positive degree works for all $i$.
\end{lemma}

All the constants in this lemma are sequence bounds.  In particular,
the phase condition is tested along arbitrary subsequences and with
arbitrarily chosen divisors.  This quantifier has two useful
consequences.  First, at a fixed growing parameter and after the
allowed scaling, there are $C>0$ and a sufficiently small fixed
$\eta>0$ such that all phases at $a_i(t_i,T)<\eta$ have their small
lifts bounded in absolute value by $Ca_i(t_i,T)$, eventually.
Otherwise one selects $a_i(t_i,T_i)\to0$ with an unbounded ratio.
Second, the condition also applies to integral multiples
$w_i=e_i\ord_{T_i}$ whenever $a_i(t_i,w_i)\to0$.  Indeed,
$a_i(t_i,T_i)\to0$, and
\begin{equation}\label{eq:B-multiple-phase}
 e_i\theta_i(T_i)=O\bigl(a_i(t_i,w_i)\bigr)\longrightarrow0,
 \qquad
 e_i\theta_i(T_i)\ge-o\bigl(a_i(t_i,w_i)\bigr).
\end{equation}
Thus multiplication introduces no wrap around the phase circle.  This
observation is needed each time a divisorial valuation restricts to a
multiple of an angular prime.

\subsection{Seifert bundles and two consequences of comparison}

We use the graded-algebra description of Seifert bundles and its
compactification \cite[Definition~5, Theorem~7 and Section~14]{KollarSeifert}.
The discrepancy and phase formulas needed here are proved below.

\begin{lemma}[Bundle construction and discrepancies]\label{lem:B-seifert}
Let $V$ be projective normal, and let $H$ be a $\Q$-Cartier rational
divisor.  There is a projective normal compactified Seifert bundle
$p:Y\to V$ with disjoint $\Q$-Cartier zero and infinity sections
$E_0,D$ and a rational fibre coordinate $z$ such that
\begin{equation}\label{eq:B-bundle-identities}
 \divisor(z)=p^*H+E_0-D,
 \qquad
 K_Y+E_0+D=p^*\bigl(K_V+R(H)\bigr).
\end{equation}
The second equality uses compatible canonical divisors.  The divisor
$D$ is relatively ample; if $H$ is nef, then $D$ is nef, and if $H$
is nef and big, then $D+p^*H$ is nef and big.

In these identities, pullback of a rational Weil divisor along the
equidimensional map $p$ is defined at the generic points of its prime
components, using the multiplicities of their inverse images.  This
agrees with Cartier pullback whenever the divisor is $\Q$-Cartier.
Orders of $H$, $E_0$, and $D$ on higher models use their Cartier
b-pullbacks.

Suppose $B\ge R(H)$ and $\mathbf M$ is a nef $\Q$-Cartier
b-divisor on $V$.  Give $Y$ boundary $p^*(B-R(H))$ and nef part
$p^*\mathbf M$.  Write $a$ for the generalized discrepancy on $V$
and $A$ for that on $Y$.  If a primitive divisorial valuation $v$ of
$k(Y)$ restricts nontrivially to $w$ on $k(V)$, then
\begin{equation}\label{eq:B-bundle-discrepancy}
 A(v)=a(w)+|h(v)|+\operatorname{def}(v),
 \qquad
 h(v)=v(E_0)-v(D)\equiv-w(H)\pmod\Z,
\end{equation}
where $\operatorname{def}(v)$ is a nonnegative integer.  It vanishes
precisely for the Gauss, or fibre-torus-invariant, valuations.
If $w=0$, the primitive valuations are the ordinary places of the
generic $\mathbb P^1$, and have discrepancy $1$.
Consequently klt base data give klt data on $Y$ without the two
sections; adding both sections gives lc data whose non-klt locus is
contained in the sections.
\end{lemma}

\begin{proof}
The two affine charts are the relative spectra of the algebras of
rational monomials $gz^j$ satisfying
$\divisor(g)+jH\ge0$, with $j\ge0$ and $j\le0$, respectively.
If $\ell H=\divisor(b)$ locally, these charts are finite cyclic
quotients of the ordinary line charts obtained using $b^{1/\ell}$
and the coordinate $z/b^{1/\ell}$.  They glue to the stated normal
projective bundle.  This description gives the two sections, their
$\Q$-Cartier property, and the first equality in
\eqref{eq:B-bundle-identities}.  The restriction of $D$ to itself is
numerically $H$, while its restriction to $E_0$ is trivial; this and
relative ampleness give the positivity assertions.

At a prime $P$ of $V$ the reduced vertical divisor has multiplicity
$r_P$.  Computing the relative canonical divisor by $d\log z$
therefore gives the second equality in
\eqref{eq:B-bundle-identities}, including exactly the coefficient
$1-1/r_P$.

For the discrepancy formula, write $w=e\ord_T$ on a smooth base
model and choose a uniformizer $x$ at $T$.  Wedge a base logarithmic
volume form with $d\log z$.  The additional logarithmic order is a
nonnegative integer; it is zero exactly when the residue of the
valuation-zero monomial $z^e/x^{v(z)}$ is transcendental over
$\kappa(T)$.  This is the Gauss condition.  Removing the two
sections from the boundary adds $v(E_0)+v(D)=|h(v)|$, since their
supports are disjoint.  Evaluating the first identity in
\eqref{eq:B-bundle-identities} gives the congruence.  The same
calculation on a determination of the nef part proves the generalized
formula.  For trivial restriction the assertion is the usual
discrepancy calculation on $\mathbb P^1_{k(V)}$.
\end{proof}

In the setting of \eqref{eq:B-family}, the data without the sections
have total logarithmic divisor
\begin{equation}\label{eq:B-bundle-total}
 K_Y+p^*(B(t)-R(H))+p^*\mathbf M(t)_Y
 \sim_{\Q}(t+1)p^*H-2D.
\end{equation}
At $t=0$, adding both sections instead gives effective generalized
lc Calabi--Yau data.

\begin{lemma}[Clearing the small denominators]\label{lem:B-denominator-clearing}
In bounded dimension, suppose $V$ is projective of Fano type and
carries effective generalized klt Calabi--Yau data with
$B\ge R(H)$.  After multiplying $H$ by a bounded positive integer,
there is an ordinary klt log Fano boundary $\Theta\ge R(H)$.
For this scaled divisor the Seifert bundle is of Fano type.
The assertion does not require a bound on the Cartier index of the
nef b-divisor.
\end{lemma}

\begin{proof}
Take a small $\Q$-factorialization.  The class $-(K_V+R(H))$ is
pseudo-effective, because it is the sum of $B-R(H)$ and the trace
of a b-nef divisor.  Moreover $(V,R(H))$ is ordinarily klt:
its discrepancies dominate the given generalized discrepancies,
by $B\ge R(H)$ and b-nef negativity.  Run its anticanonical MMP,
using Fano type to obtain a nef model.  Ordinary lc comparison
holds along this MMP by the negativity argument of
Section~\ref{sec:A}.  The coefficients of
$R(H)$ belong to the fixed hyperstandard set
$\{1-1/r:r\in\N\}$.  The bounded ordinary lc complement theorem
\cite[Theorems~1.7--1.8]{BirkarComplements}, followed by pullback
through this MMP, gives an ordinary lc boundary
$\Delta\ge R(H)$ with $n(K_V+\Delta)\sim0$, where $n$ is bounded
in terms of dimension.  The pullback dominates $R(H)$ because the
MMP has the anticanonical sign; this is the same complement
pullback used in Lemma~\ref{lem:A-lc-place-tools}.

If $r_P>n$, the $1/n$-quantization of $\Delta$ forces
$\operatorname{coeff}_P(\Delta)=1$.  Multiply $H$ by the least
common multiple of the integers at most the bound for $n$.
All smaller denominators disappear.  Every remaining component of
$R(H)$ is therefore contained in $\lfloor\Delta\rfloor$.
Let $\Theta_0$ be any klt log Fano boundary on $V$.  For a
sufficiently small positive rational $\rho$, chosen separately for
each variety,
\[
 \Theta=(1-\rho)\Delta+\rho\Theta_0
\]
is klt, dominates the new $R(H)$, and has ample negative log
canonical class.  No uniform lower bound for $\rho$ is needed.

On the bundle use boundary
$p^*(\Theta-R(H))+(1-\rho')(E_0+D)$, with $\rho'>0$ sufficiently
small.  Lemma~\ref{lem:B-seifert} proves klt.  If
$A=-(K_V+\Theta)$, its negative log canonical class is
$p^*(A-\rho'H)+2\rho'D$, which is ample for small $\rho'$ by
relative ampleness and base positivity.  Thus the bundle is of
Fano type.
\end{proof}

\begin{lemma}[An unbounded-index obstruction]\label{lem:B-unbounded-index}
Let $F_i$ be projective of positive bounded dimension and let $G_i$
be ample $\Q$-Cartier rational divisors such that $(F_i,R(G_i))$
is klt and
\[
 -(K_{F_i}+R(G_i))\sim_{\Q}I_iG_i,
 \qquad I_i>0,\quad I_i\longrightarrow\infty.
\]
Let $C_i$ be the compactified Seifert cone obtained by contracting
the zero section, and denote its infinity divisor by $D_i$.
Then $C_i$ is klt of Fano type, $D_i$ is ample, and
$-K_{C_i}\sim_{\Q}(I_i+1)D_i$.  It is impossible that
\begin{equation}\label{eq:B-cone-small-height}
 v_i(D_i)=o\bigl(A_{C_i}(v_i)\bigr)
\end{equation}
for every sequence of primitive divisorial valuations with
$A_{C_i}(v_i)\to0$.
\end{lemma}

\begin{proof}
The contraction exists by semiampleness of the infinity divisor.
One can also construct it by the finite fibre-power map to the
ordinary compactified bundle for a very ample Cartier multiple of
$G_i$, contracting its zero section and then Stein factorizing.
The bundle identity gives
\[
 K_Y+(1-I_i)E_0+(1+I_i)D\sim_{\Q}0.
\]
The pair with subboundary $(1-I_i)E_0$ is sub-klt by
\eqref{eq:B-bundle-discrepancy}; the discrepancy of the contracted
zero section is $I_i>0$.  This proves the assertions about the
cone and its canonical divisor.  Near infinity the discrepancy
formula reads, for $u=v(D)>0$,
\begin{equation}\label{eq:B-infinity-formula}
 A_C(v)=A_{F,R(G)}(w)+u+\operatorname{def}(v),
 \qquad u\equiv w(G)\pmod\Z.
\end{equation}

If \eqref{eq:B-cone-small-height} held, apply
Lemma~\ref{lem:A-comparison} over a point to the nonincreasing
family with discrepancies $A_C-tv(D)$, boundary $tD$, and nef part
$(I_i+1-t)D$ on a range tending to infinity.  The phase-free
ratio hypothesis is exactly \eqref{eq:B-cone-small-height}.
The comparison lemma would make this family klt on every fixed
bounded parameter range, whereas at $t=1$ the prime $D$ has
discrepancy zero.  This is a contradiction.
\end{proof}

\begin{lemma}[Bounding a principalization order]\label{lem:B-principal-order}
Let $F_i$ be projective of Fano type in bounded dimension, and let
$P_i\sim_{\Q}0$ be $\Q$-Cartier rational divisors.  Suppose that
$F_i$ carries effective generalized klt Calabi--Yau data with
discrepancies $a_i$ and boundary at least $R(P_i)$.  Suppose,
after a fixed integral scaling if necessary, that the signed phases
of $P_i$ at every sequence with $a_i(T_i)\to0$ lie in
$[-o(a_i(T_i)),O(a_i(T_i))]$.  Then the least positive integer $q_i$
for which $q_iP_i$ is principal is bounded.

More generally, the same cyclic-cover argument applies whenever
the induced isotrivial fibration below has
$u=o(A)$ at its small-discrepancy divisors over zero.  This
formulation permits the argument to be used after an independent
estimate of the height $u$.
\end{lemma}

\begin{proof}
Apply Lemma~\ref{lem:B-denominator-clearing} to obtain, after a
further bounded scaling, ordinary klt log Fano data dominating
$R(P_i)$.  This does not affect boundedness of the original
principalization orders.  A small $\Q$-factorialization preserves
all the data and these orders.  Suppress the index and suppose
for contradiction that $q\to\infty$.  Write
\[
 \divisor(\phi)=qP,
 \qquad
 \widetilde F=\operatorname{Norm}_F\bigl(k(F)(r)\bigr),
 \qquad r^q=\phi.
\]
The degree is $q$ by minimality of $q$ and the Kummer criterion
over a field containing all roots of unity.  The ramification in
codimension one is encoded by $R(P)$.  In particular the log
pullbacks of both the ordinary log Fano boundary and the
generalized boundary are effective and klt.

Form the diagonal quotient
\begin{equation}\label{eq:B-cyclic-fibration}
 X=(\widetilde F\times\mathbb P^1_y)/\mu_q
 \longrightarrow\mathbb P^1_{y^q},
\end{equation}
where $\mu_q$ acts inversely on $y$ and $r$, so that $yr$ is
invariant.  The map from the product to $X$ is quasi-\'etale:
a nonidentity element has a fixed locus of codimension at least
one on $\widetilde F$ and only the two fixed points on
$\mathbb P^1$.  The generic fibre of \eqref{eq:B-cyclic-fibration}
is connected.  Log pulling the ordinary log Fano data to the
product and descending proves that $X$ is of Fano type over the
curve.

Pull back the generalized data to the product and descend them.
They give effective generalized klt Calabi--Yau data over the
curve.  The invariant nef b-divisor descends on a sufficiently
high finite quotient; its $\Q$-Cartier property and relative
rational triviality descend by taking norms.  For a primitive
valuation over zero put
$u=v(y^q)/q>0$, and let $w$ be its restriction to $k(F)$.
The product discrepancy calculation gives
\begin{equation}\label{eq:B-cyclic-height}
 A(v)\ge a(w)+u,
 \qquad u\equiv-w(P)\pmod\Z,
\end{equation}
where $a(w)=0$ if $w=0$.  The congruence follows by evaluating
the invariant function $yr$.  These formulas can equally be
computed upstairs and divided by the ramification index, using
log pullback and homogeneous discrepancies.

If $A(v_i)\to0$, trivial restriction is impossible because it
would make $u_i$ a positive integer.  For nontrivial restriction,
\eqref{eq:B-multiple-phase} and \eqref{eq:B-cyclic-height} give
$u_i=o(A(v_i))$: all terms tend to zero, so there is no wrap,
and the lower bound for the signed phase bounds the positive
number $u_i$ from above.

Apply Lemma~\ref{lem:A-comparison} over the curve to
$A-t\,v(f^*[0])/q$.  These are effective generalized
Calabi--Yau data obtained by adding $tf^*[0]/q$ to the boundary;
the nef part is unchanged.  The parameter range may tend to
infinity.  The reduced zero fibre has discrepancy $1$ and
multiplicity $q$, hence its initial discriminant discrepancy
satisfies $0<b_0(0)\le1/q$.  Since the perturbation is a pure
base pullback,
\[
 b_t(0)=b_0(0)-t/q,
 \qquad b_1(0)\le0.
\]
This contradicts the ratio conclusion of
Lemma~\ref{lem:A-comparison}, because $b_0(0)\to0$.
\end{proof}

\subsection{Proof of phase nonvanishing}

We now prove the main nonvanishing lemma.  The first task is to subtract
small positive phase corrections from the polarization and show that the
remaining divisor is not big.  Its subsequent MMP will either contradict
the unbounded-index obstruction or produce a lower-dimensional base to
which the same phase hypotheses descend.

\begin{proof}[Proof of Lemma~\ref{lem:B-phase-nonvanishing}]
\emph{Choosing the phase corrections.}
We argue by dimension induction.  The assertion for a point is
immediate.  If the assertion fails in positive dimension, pass to
a sequence and choose integers $N_i\to\infty$ such that
\begin{equation}\label{eq:B-no-sections}
 |\lfloor mH_i\rfloor|=\varnothing
 \quad\text{for all integers }1\le m\le N_i.
\end{equation}
Choose rational $k_i\to\infty$ sufficiently slowly relative to
$N_i$ and $L_i$.  Use the phase hypothesis at $k_i$ and
Lemma~\ref{lem:B-denominator-clearing}, pass to a subsequence,
and scale $H_i$ by fixed positive integers.  Reparametrize the
given families by these same factors, and decrease $N_i$ by
bounded factors.  We retain \eqref{eq:B-no-sections}, obtain the
phase hypothesis without a displayed scaling at $k_i$, and have
Fano type Seifert bundles.  These bounded changes never replace
the given discrepancy families by different ones.

Fix a sufficiently small constant $\eta>0$.  Extract precisely
the exceptional primes $T$ with $a(k_i,T)<\eta$ whose small signed
phase $f_T$ of $H(T)$ is positive.  Include the original primes
with the same property in the list of corrections.  This is a
finite list: monotonicity gives $a(0,T)\le a(k_i,T)<\eta$, and
the extraction and finiteness argument of
Section~\ref{sec:A} applies to the klt data $a(0)$.
Write the resulting Fano type model as $U\to V$.  For small fixed
$\eta$ we have, eventually,
\begin{equation}\label{eq:B-positive-corrections}
 0<f_T\le C a(k_i,T)\le C\eta,
 \qquad
 C_U=\sum_T f_TT,
 \qquad P_U=H_U-C_U,
\end{equation}
where $H_U$ is the pullback of $H$.
We first prove that $P_U$ is not big, eventually.

\medskip
\noindent\emph{The threshold contradiction under bigness.}

Suppose instead that $P_U$ is big.  On the Seifert bundle let $A$
be the discrepancy of the data at $k_i$ without the sections.
Choose a small fixed $0<\delta<1/4$, sufficiently small relative
to $\eta$ and the phase bound.  Extract every exceptional prime
with $A<\delta$, and let $\mathcal S$ be the reduced sum of
\emph{all} primes with $A<\delta$ on this extraction, including
original primes.  Lemma~\ref{lem:B-seifert} shows that they are
Gauss valuations.  Their extraction is funded by the effective
zero-parameter generalized Calabi--Yau data with both sections,
mixed with a sufficiently small ordinary klt log Fano
contribution as in Section~\ref{sec:A}.  Thus the extraction is
of Fano type, and the trace boundary of $A$ is effective on it.

Put $\mathbf H=p^*H$, and denote by $\mathbf D$ the b-pullback
of $D$.  The b-divisor
\[
 \mathbf J=\mathbf D+\mathbf H
\]
is nef and big and has integral Weil trace on the original
bundle: the multiplicities $r_P$ clear the coefficients of
$H$ at all vertical primes.  Let $\tau$ be the pseudo-effective
threshold for subtracting $\mathcal S$ from its trace on the
extraction.  We claim that
\begin{equation}\label{eq:B-threshold-small}
 \mathcal S\ne0,
 \qquad k_i\tau_i\longrightarrow0.
\end{equation}

To prove the claim, if it fails pass to a subsequence and choose
integers
$k_i,\tau_i^{-1}\ll m_i\ll N_i$; omit $\tau_i^{-1}$ when
$\mathcal S=0$.  Round the coefficients on the extracted primes
to obtain
\[
 Q=\sum_{T\subset\mathcal S}q_TT,
 \quad 1\le q_T\le2,
 \qquad N=m\mathbf J_U-Q
\]
big and integral.  All nonintegral coefficients of $m\mathbf J_U$
are on the extracted primes, so this rounding suffices.  When
$\mathcal S=0$, use $N=mJ$.  Run the $N$-MMP to a nef and big
model.  Here and in what follows a prime denotes the output
trace, and a star denotes its Cartier b-pullback to a common
resolution.  Set
\[
 e_H=(H')^*-\mathbf H,\qquad
 e_D=(D')^*-\mathbf D,\qquad e_J=e_H+e_D.
\]
B-nef negativity and MMP negativity give
\begin{equation}\label{eq:B-rounded-errors}
 e_H,e_D\ge0,
 \qquad
 m e_J\le(Q')^*-Q^*,
 \qquad (Q')^*\le2(\mathcal S')^*.
\end{equation}
The total logarithmic divisor for $A$ has trace-pull error
$(k_i+1)e_H-2e_D$, by \eqref{eq:B-bundle-total}.  Its effective
trace boundary dominates $(1-\delta)\mathcal S'$.  Comparing
generalized and ordinary discrepancies on the output therefore
gives
\begin{align}
 A_{\mathrm{ord}}'
 &\ge A+(1-\delta)(\mathcal S')^*-(k_i+1)e_H\notag\\
 &\ge A+
 \left(1-\delta-\frac{2(k_i+1)}m\right)(\mathcal S')^*
 \ge A+\frac12(\mathcal S')^* .
 \label{eq:B-ordinary-gap}
\end{align}
The comparison also discards only a nonnegative trace-pull error
of the nef part.

This gives a gap at contracted primes as well as surviving ones.
Indeed, every primitive $T$ with $A(T)<\delta$ was put on the
extraction, and hence $Q^*(T)=q_T\ge1$, even if the subsequent
MMP contracts $T$.  Equation~\eqref{eq:B-rounded-errors} implies
\begin{equation}\label{eq:B-contracted-prime-bound}
 (Q')^*(T)\ge Q^*(T)\ge1,
 \qquad (\mathcal S')^*(T)\ge\frac12.
\end{equation}
Such a prime has $A_{\mathrm{ord}}'(T)\ge1/4$ by
\eqref{eq:B-ordinary-gap}; all other primitive primes have
$A_{\mathrm{ord}}'(T)\ge A(T)\ge\delta$.  Thus the output is
ordinarily $\min\{\delta,1/4\}$-lc.

Fano type makes $-K'$ pseudo-effective, so
$N'-K'$ is pseudo-effective.  Integral-Weil effective
birationality \cite[Theorem~1.1]{BirkarPolarised} supplies a
nonzero section of a bounded multiple of $N'$.  MMP negativity
pulls it back to the same multiple of $N$, and then to a bounded
multiple of $m\mathbf J$.
The fibre torus acts on the latter section space, so it contains
a weight vector $gz^j$.  A section of $\ell J$ satisfies
\begin{equation}\label{eq:B-angular-section}
 0\le j\le\ell,
 \qquad \divisor(g)+(j+\ell)H\ge0.
\end{equation}
Indeed the coefficients of its divisor along $E_0,D$ are
$j,\ell-j$, respectively, by \eqref{eq:B-bundle-identities}.
Thus it gives a section of $\lfloor (j+\ell)H\rfloor$ in positive
degree at most $2\ell=O(m_i)\ll N_i$, contrary to
\eqref{eq:B-no-sections}.  This proves
\eqref{eq:B-threshold-small}.

Run the MMP of $J-\tau\mathcal S$ to its semiample nonbig model
and let its Iitaka fibration be $Y'\to Z_\tau$.  The threshold
is rational because the extraction is of Fano type.  Negativity
now gives
\begin{equation}\label{eq:B-threshold-errors}
 0\le e_J\le\tau\bigl((\mathcal S')^*-\mathcal S^*\bigr).
\end{equation}
Choose a small fixed $c$ with $0<c<1-\delta$.  On this output
use the boundary $B_Y'-c\mathcal S'$ and nef b-divisor
\begin{equation}\label{eq:B-threshold-nef-part}
 \mathbf M^\#
 =\mathbf M(k_i)+
 \left(2+\frac c\tau\right)\mathbf D+
 \left(\frac c\tau-k_i-1\right)\mathbf H.
\end{equation}
This boundary is effective.  Both displayed coefficients in the
nef part are eventually nonnegative by
\eqref{eq:B-threshold-small}; in particular the nefness is
absolute, and uses no Cartier-index bound.  Its total
logarithmic divisor is $\Q$-linearly equivalent to
$(c/\tau)(J'-\tau\mathcal S')$.  The discrepancies are
\begin{equation}\label{eq:B-modified-gap}
 A^\#=A+c\left((\mathcal S')^*-\frac{e_J}{\tau}\right)
 \ge A+c\mathcal S^*.
\end{equation}
Every primitive prime with $A<\delta$ has $\mathcal S^*(T)=1$,
including a prime contracted along the MMP.  The remaining
primes have $A\ge\delta$.  Thus $A^\#$ has a uniform generalized
gap $\min\{c,\delta\}$.

On a geometric generic fibre $F$ of $Y'\to Z_\tau$ these are
effective generalized Calabi--Yau data.  Restriction can be
computed using horizontal divisors on all higher models over
the generic field, and then separable extension of constants;
it is crepant and preserves the displayed gap.  The fibre is
of Fano type and rationally connected.  Apply
\cite[Theorem~1.7]{BirkarCY} to $A^\#$ on $F$ and take a bounded
normal model isomorphic to it in codimension one.  Normalization
of the bounded models is still bounded after stratification and
spreading, and does not add divisors to a codimension-one
isomorphism.

The effective boundary $B_Y'-c\mathcal S'$ has coefficients at
least $1-\delta-c$ on $\mathcal S'$.  Intersecting its
Calabi--Yau equality with a fixed bounded polarization bounds
its degree: the trace of the nef b-part is pseudo-effective,
and the canonical degree is bounded in the bounded family.
Consequently
\begin{equation}\label{eq:B-small-degrees}
 \deg\mathcal S'=O(1),\qquad
 \deg J'=O(\tau),\qquad
 \deg\mathbf H_F,\ \deg\mathbf D_F=O(\tau).
\end{equation}
Here $J'\sim_{\Q}\tau\mathcal S'$ on the fibre, and the last
two degrees are nonnegative and sum to that of $J'$.

We now use the assumed bigness of $P_U$.  It gives an effective
rational divisor $\Sigma\sim_{\Q}H$ such that for every $v$
listed in $\mathcal S$, with restriction $w$, one has
\begin{equation}\label{eq:B-strict-improvement}
 w(\Sigma)>v(D)-v(E_0).
\end{equation}
For completeness, if $v(D)>0$, write $w=e\ord_T$.
The discrepancy formula gives $ea(k_i,T)<\delta$ and
$v(D)<\delta$.  Since $e|f_T|\le Cea(k_i,T)$ and $\delta$ is
sufficiently small, the congruence has no wrap and
$v(D)=ef_T$.  Thus $f_T>0$, and this prime was corrected in
$P_U$.  An effective representative of the big class $P_U$
may be chosen to vanish strictly at these primes and at any of
the remaining finitely many projected centres.  Every listed
valuation has nontrivial angular restriction, since trivial
restriction would give $A=1$; its centre on $U$ is therefore proper.
Reserve a tiny ample summand and choose a sufficiently divisible
member through all these centres.  Adding back $C_U$ gives
\eqref{eq:B-strict-improvement}.  Pushing down yields an
effective representative of $H$ on $V$; the representative on
$U$ is its pullback, since it differs from $H_U$ by a rational
principal divisor.

Every homogeneous normalized threshold section has divisor
\begin{equation}\label{eq:B-threshold-section}
 (1-b)D+bE_0+p^*\Xi,
 \qquad 0\le b\le1,
 \qquad \Xi\sim_{\Q}(1+b)H.
\end{equation}
If $b<1$, replace $b$ by $b+\alpha$ and $\Xi$ by
$\Xi+\alpha\Sigma$ for a sufficiently small positive rational
$\alpha$.  The divisor remains effective, while every
$\tau\mathcal S$ constraint improves strictly by
\eqref{eq:B-strict-improvement}.  The finite list of strict
improvements would increase the pseudo-effective threshold,
a contradiction.  Hence $b=1$ for every such section.
All these systems are torus invariant, because the listed
valuations and the original $J$ are torus invariant.  Equal-degree
ratios therefore lie in $k(V)$, by their weight decompositions
and $b=1$.  The valuation
$v_0=\ord_{E_0}$ is trivial on this field, so it is horizontal
over the Iitaka base $Z_\tau$.

We obtain a contradiction on the geometric generic fibre.
Tailor a bounded ordinary lc complement preserving $v_0$ as
an lc place, using the effective zero-parameter generalized
Calabi--Yau data with the two sections and the Fano type
funding of Lemma~\ref{lem:A-lc-place-tools}.  This may first
be done on a small $\Q$-factorialization of the fibre, using
the horizontal lift of $v_0$, and then transferred to the
bounded model.  Work on its homogeneous horizontal chart.
Proposition~\ref{prop:A-uniform-charts} bounds the Lipschitz
constant of $A^\#$.  The difference $A^\#-A$ is the order of
\[
 \frac c\tau\left(\mathbf J-(J'-\tau\mathcal S')^*\right).
\]
The parenthesized b-divisor has effective trace
$\tau\mathcal S'$ of degree $O(\tau)$ on this fibre and is
b-nef up to a rational principal divisor.  The shrinking-degree
estimate of Corollary~\ref{cor:A-horizontal-variants} therefore
gives an $O(1)$ Lipschitz bound for the difference, and hence
for $A$.

The same estimate gives an $O(\tau)$ Lipschitz bound for
$h_0(v)=v(E_0)$.  Indeed $E_0$ has effective trace and is
b-linearly equivalent to $\mathbf D-\mathbf H$, whose two nef
terms have degree $O(\tau)$ by \eqref{eq:B-small-degrees}.
In applying the estimate one replaces these nef terms by
effective representatives with arbitrarily small error,
avoiding the two centres in question.  The function-degree
argument on the bounded chart then applies: a rational
function with pole degree strictly below $1$ is constant.
This is why the Lipschitz bound tends to zero, rather than
merely remaining bounded.

Let $w_0$ be a chart vector representing $v_0$.  We have
$h_0(w_0)=A(w_0)>0$.  The shrinking Lipschitz bound implies
$|w_0/A(w_0)|\to\infty$.  The horizontal approximation in
Lemma~\ref{lem:A-lattice-approximation}, applied to this ray,
gives integral divisorial vectors $n$ and numbers $s\to0^+$
such that
\[
 A(n)=O(s)\longrightarrow0,
 \qquad h_0(n)\ge s/2.
\]
On the other hand, any such small-discrepancy valuation with
$h_0>0$ has $v(D)=0$.  Formula~\eqref{eq:B-bundle-discrepancy}
and the lower signed phase bound at $a(k_i)$, including
\eqref{eq:B-multiple-phase}, imply $h_0=o(A)$; trivial
restriction would instead make $h_0$ a positive integer and
is impossible.  This contradicts the two displayed estimates.
We have proved that $P_U$ is not big.

\medskip
\noindent\emph{Running the correction MMP.}

Diagonalize the preceding conclusion to choose
$\eta_i\to0^+$ sufficiently slowly so that $P_U$ is still
not big.  Run its MMP.  All parameters used from now on may
grow more slowly, remain below $k_i$, and be $o(1/\eta_i)$.
On the output put
$\mathcal T'$ equal to the reduced sum of surviving correction
primes.  With $\mathbf H$ now denoting the original angular
b-pullback, MMP negativity and b-nefness give
\begin{equation}\label{eq:B-correction-errors}
 0\le C_U^*\le
 E:=\mathbf H-(P')^*
 \le(C')^*,
 \qquad C'\le O(\eta_i)\mathcal T'.
\end{equation}
More explicitly, if $e_H=(H')^*-\mathbf H\ge0$, then
$E=(C')^*-e_H$, and $P$-negativity gives
$(C')^*-C_U^*\ge e_H$.

\medskip
\noindent\emph{Excluding a Mori fibre output.}
Suppose first that $P_U$ is not pseudo-effective.  On the
geometric generic fibre $F$ of the final anti-$P'$ Mori
contraction, put $G=-P'$, which is ample.  Choose a new slow
parameter $s_i\to\infty$, with $s_i\eta_i\to0$, and obtain
the phase condition at $a(s_i)$ after a fixed scaling.
Scale $H,P_U,C_U$ together.  The correction coefficients
remain $O(\eta_i)$, and $P_U$ is integral at correction
primes: subtracting their phases was precisely what made
those coefficients integral.  Thus $R(G)$ is supported
away from $\mathcal T'$, and the pushed boundary satisfies
\[
 B(s_i)'\ge R(G)+(1-o(1))\mathcal T'.
\]
Comparing with ordinary discrepancies on $F$, the only
adverse total-divisor error is $s_ie_H$, which is at most
$s_i(C')^*$.  B-nef negativity for the moduli part and
\eqref{eq:B-correction-errors} therefore give
\begin{equation}\label{eq:B-negative-discrepancy}
 A_{F,R(G)}
 \ge a(s_i)+\tfrac12(\mathcal T')^*.
\end{equation}
In particular $(F,R(G))$ is klt.  A uniform ordinary gap is
not asserted or needed in this case.

The relative Picard number of the Mori contraction is one.
The classes used here, including $R(G)$, are restrictions
from this model.  Horizontal prime coefficients restrict
unchanged over the generic field and split with unchanged
multiplicities after a separable constant extension, so the
same class comparisons hold on the geometric generic fibre.
There we have
\[
 G+H'=C'\le O(\eta_i)\mathcal T',
\]
where $H'$ is pseudo-effective.  The zero-parameter
generalized Calabi--Yau equality also shows that
$-(K_F+R(G))$ dominates a fixed positive multiple of
$\mathcal T'$ in this one-dimensional space of restricted
classes: its boundary contributes $(1-o(1))\mathcal T'$,
and the moduli trace is pseudo-effective.  Consequently
\begin{equation}\label{eq:B-large-index}
 -(K_F+R(G))\sim_{\Q}I_iG,
 \qquad I_i\longrightarrow\infty.
\end{equation}
Initially this is a numerical proportionality, with
$I_i\gg1/\eta_i$.  It is rational linear proportionality
after clearing denominators because $F$ is rationally
connected: pull numerically trivial Cartier multiples to a
smooth rationally connected resolution, where they are
torsion, and descend the resulting rational equivalence.

Consider the compactified cone for $(F,G)$ and a primitive
divisorial valuation with ordinary discrepancy $A\to0$ and
infinity height $u>0$.  Its angular restriction $w$ has
$A_{F,R(G)}(w)\le A$ and $u\le A$ by
\eqref{eq:B-infinity-formula}.  Moreover
\eqref{eq:B-correction-errors} and
\eqref{eq:B-negative-discrepancy} imply
\begin{equation}\label{eq:B-negative-error-small}
 0\le E(w)\le O(\eta_i)A_{F,R(G)}(w)\le O(\eta_i)A.
\end{equation}
Divisorial tests on the geometric fibre lift to homogeneous
horizontal tests on the original model.  Thus the signed
phase bound for $H$ applies, and
\[
 u\equiv w(G)\equiv-w(H)+E(w)\pmod\Z.
\]
All three quantities have small signed representatives.
There is no wrap, so the lower phase bound and
\eqref{eq:B-negative-error-small} give $u=o(A)$.
If $w=0$, the congruence makes $u$ a positive integer,
already impossible.  Valuations with $u=0$ satisfy the
same conclusion trivially.  This contradicts
Lemma~\ref{lem:B-unbounded-index} and
\eqref{eq:B-large-index}.

\medskip
\noindent\emph{Descent from the semiample output.}
We are left with a semiample nonbig $P'$ and its contraction
$g:U'\to W$.  For a newly trimmed parameter range tending
to infinity, with $t\eta_i\to0$ uniformly on that range,
use boundary and nef part
\begin{equation}\label{eq:B-new-family}
 \widetilde B(t)=B(t/2)'-tC',
 \qquad
 \widetilde{\mathbf M}(t)=\mathbf M(t/2)+(t/2)\mathbf H.
\end{equation}
Their total logarithmic divisor is $tP'$ up to rational
linear equivalence, and their discrepancies are exactly
\begin{equation}\label{eq:B-new-discrepancy}
 \widetilde a(t)=a(t/2)+tE.
\end{equation}
To see the equality, first pushing the old total divisor
contributes $-(t/2)e_H$; adding the nef term and subtracting
$tC'$ contributes another $-(t/2)e_H+t(C')^*$.
Their sum is $tE$.

The new boundary dominates $R(P')$.  At correction primes
$P'$ is integral, the old boundary has coefficient
$1-o(1)$, and the amount subtracted is $O(t\eta_i)=o(1)$.
Away from them the old denominator condition is unchanged.
The nef part is b-nef; \eqref{eq:B-new-discrepancy} and
$E\ge0$ show that the discrepancies are increasing and
klt, with effective Calabi--Yau data at zero.

This family has the required signed phase condition for
$P'$.  Indeed, at a new growing parameter, apply the
original condition to $a(t/2)$ and scale all the relevant
divisors by the resulting fixed integer.  If
$\widetilde a(t,T)\to0$, then
\[
 a(t/2,T)\le\widetilde a(t,T),
 \qquad E(T)\le\widetilde a(t,T)/t
      =o(\widetilde a(t,T)).
\]
Since $P'(T)=H(T)-E(T)$, subtracting this negligible
nonnegative order preserves the lower $-o(\widetilde a)$
bound and the upper $O(\widetilde a)$ bound for the phase.

On the geometric generic fibre of $g$, the divisor $P'$
is $\Q$-linearly trivial.  Its integral principalization
order is bounded by Lemma~\ref{lem:B-principal-order}:
along any putative unbounded-order subsequence choose one
growing parameter with the just-proved phase bound, use
the restricted effective generalized klt Calabi--Yau
data, and clear the remaining bounded small denominators
as in that lemma.  These operations use only fixed
multipliers.  A small $\Q$-factorialization of the
geometric fibre preserves all the divisor orders and
permits the cyclic-cover construction.

The principalization descends from the geometric generic
fibre to its field of definition.  Indeed being Cartier
descends faithfully flatly, and the geometrically trivial
line bundle is trivial by Hilbert 90; equivalently its
one-dimensional space of trivializing rational functions
descends.  Thus for a bounded positive integer $q$ and
some $\phi\in k(U')^*$, the divisor
$qP'-\divisor(\phi)$ has no horizontal terms.  There is
then an exact equality
\begin{equation}\label{eq:B-exact-descent}
 qP'-\divisor(\phi)=g^*H_W
\end{equation}
for an ample $\Q$-Cartier rational divisor $H_W$.
To justify exactness, semiampleness already gives
$P'\sim_{\Q}g^*A_W$ with $A_W$ ample.  After clearing
denominators the difference between the left side of
\eqref{eq:B-exact-descent} and $qg^*A_W$ is principal
and vertical.  A rational function with vertical divisor
restricts to a nowhere vanishing regular function on
the projective generic fibre, hence belongs to $k(W)^*$
because $g$ is a contraction.  Absorb this rational
principal divisor into $qA_W$ to obtain $H_W$.

If $W$ is a point, \eqref{eq:B-exact-descent} and $E\ge0$
already give a section of a bounded multiple of $H$.
Otherwise $W$ has strictly smaller dimension than $V$.
It is of Fano type.  Apply the generalized canonical
bundle formula to \eqref{eq:B-new-family}, and
replace the old parameter $s$ by $s=qt$.  The induced
data on $W$ are effective generalized data with total
$tH_W$, klt at zero and with nondecreasing discrepancies.
Their nef parts are b-nef and $\Q$-Cartier.

We check the denominator condition on $W$ explicitly.
For a prime $P$ on $W$, choose a nonexceptional vertical
component $T$ above it with multiplicity $e$.  Let $r$
be the denominator of $\operatorname{coeff}_T(P')$ and
let $r'$ be that of
$\operatorname{coeff}_T(qP'-\divisor(\phi))$.
Then $r'\mid r$.  The denominator $r_P$ of $H_W$
divides $er'$, by \eqref{eq:B-exact-descent}.
Since the upstairs boundary dominates $R(P')$, the
discriminant discrepancy at $P$ satisfies
\[
 b_t(P)\le\frac1{er}\le\frac1{er'}\le\frac1{r_P}.
\]
Thus the discriminant boundary dominates $R(H_W)$.

For each growing parameter, the pure pullback phase
condition descends by Lemma~\ref{lem:A-pullback-phases},
after its allowed fixed multiple.  Every hypothesis of
Lemma~\ref{lem:B-phase-nonvanishing} is now satisfied in
smaller dimension.  Induction gives bounded-degree
sections of $H_W$.  They pull back through
\eqref{eq:B-exact-descent} to bounded-degree sections
of $P'$, and the b-divisor inequality
$\mathbf H\ge(P')^*$ in
\eqref{eq:B-correction-errors} carries them to sections
of the same bounded multiples of $H$ on $V$.
This contradicts \eqref{eq:B-no-sections} and completes
the induction.

The contradiction excludes every sequence of failures
with $N_i\to\infty$, and hence proves boundedness for
the original sequence.  Any finite initial segment
omitted in an eventual assertion can also be included:
each of its big divisors has a section in some positive
degree.  Taking a common multiple of the bounded
degrees gives the common degree in the statement.
\end{proof}
